% ================================================================
\renewcommand{\realmaccent}{columnBody}  % ruby red
\chapter{Body}
% ================================================================

% Full-width summary table, placed immediately after the chapter title so
% it shares the chapter's own opening page. Uses cuted's \strip
% environment, NOT table*/\twocolumn[...] --- both of those are floats
% (or trigger the float queue) and LaTeX will not let a floated top-block
% claim the top of a page that a sectioning command already opened, no
% matter how little else is on that page. \strip renders inline, exactly
% where it appears in the source, with no deferral.
% REAL RULE, 2026-09-22: rank 1 (D tier) is defined by being totally
% basic --- ordinary people have no Body Tags at all. Numbers rebalanced
% to weight Body toward Stat Points over tags; see the Designer's Note
% in Skills for the balancing principle shared between the two chapters.
% REAL RULE, 2026-09-23: Stat Points cut from 2/4/6/9 to 0/2/4/7 (a flat
% -2 shift) alongside cuts to Skills Known and Stuff Points --- smaller
% pools across the board, aimed at fewer dice to track and smaller rolls
% at the table, per playtesting. Body Tags unchanged.
% REAL RULE, 2026-09-24: Body no longer grants Stat Points at all. Body
% rank sets your Body Stat directly (1/2/3/5); Skills grants the points
% for Mind and Face (see Skills). With one letter per column, no
% character can be A in both, so two Stats at 5 is impossible --- the
% best spread anyone can reach is 5/3/3 (or 5/4/2). Rank 4 jumps
% straight to 5 on purpose: a monstrous body should leap, not creep.
\begin{strip}
  \begingroup\centering
  {\hdfont\Large\bfseries\color{\realmaccent} Creating a Body\par}
  \vspace{10pt}
  \endgroup
  \begingroup
  \setlength{\tabcolsep}{7pt}\renewcommand{\arraystretch}{1.35}
  \centering\begin{tabularx}{\textwidth}{@{}>{\centering\arraybackslash}p{1.6cm} Y >{\centering\arraybackslash}p{2.1cm} >{\centering\arraybackslash}p{1.6cm} >{\centering\arraybackslash}p{1.6cm}@{}}
  \rowcolor{\realmaccent}
  {\leavevmode\hdfont\bfseries\color{white}Rank} &
  {\leavevmode\hdfont\bfseries\color{white}You are\ldots} &
  {\leavevmode\hdfont\bfseries\color{white}Body Stat} &
  {\leavevmode\hdfont\bfseries\color{white}Body Tags} &
  {\leavevmode\hdfont\bfseries\color{white}Special Traits} \\
  \rowcolor{atrBoxbg}
  \rating{1} & An ordinary person --- nothing about you stands out.
    & \dice{1} & \dice{0} & \dice{0} \\
  \rowcolor{white}
  \rating{2} & You stand out --- strong, quick, or striking.
    & \dice{2} & \dice{1} & \dice{0} \\
  \rowcolor{atrBoxbg}
  \rating{3} & Exceptional --- elven grace, ogrish might, uncanny beauty.
    & \dice{3} & \dice{2} & \dice{1} \\
  \rowcolor{white}
  \rating{4} & Truly strange or monstrous --- a dragon, an angel.
    & \dice{5} & \dice{4} & \dice{2} \\
  \end{tabularx}
  \endgroup
\end{strip}

% TODO: opening fiction for Body.
\begin{readaloud}
\lipsum[11][1-3]
\end{readaloud}

\chapterart{col-body.png}

\lettrine{B}{ody} is your physical self --- the shape and substance you
carry from world to world. Your Body rank sets three things at once: your
\keyword{Body} Stat, how many
\keyword{Body Tags} describe your physical form, and --- if you rank high
enough --- which \keyword{Special Traits} let you do things no ordinary
roll of the dice can.

% ---- Section 1: Character Creation ---------------------------------
\section{At Character Creation}

Your Body rank --- set when you assigned priorities at character
creation, see the table above --- grants three things: your Body Stat,
Body Tags, and, at rank \rating{3} and \rating{4}, Special Traits. The
next three subsections cover each in turn.

\subsection*{Your Body Stat}
Your Body rank sets your \keyword{Body} Stat: \dice{1} at \rating{1},
\dice{2} at \rating{2}, \dice{3} at \rating{3}, and \dice{5} at
\rating{4}. A truly monstrous body doesn't creep upward --- it leaps.

\subsection*{Choosing Your Body Tags}
A \keyword{Body Tag} is a word in an oval that describes something your
body can simply \emph{do} --- \emph{Strong}, \emph{Fast}, \emph{Keen
Senses}. Whenever a tag clearly applies to what you're doing, add one die
for it to your pool, exactly like a Magic tag. Pick your tags from the
Common Body Tags list at the end of this chapter, or invent your own with
your GM's approval.

At Body \rating{1} you have none. That absence is the point: an ordinary
person's body doesn't do anything a Body Tag would describe --- that's
what makes them ordinary.

Body Tags are unlike every other resource on your sheet: they are
\emph{inherent}. No one can steal them, break them, or take them away --
they travel with you always. When you cross into a new realm, your tags
stay exactly as powerful, even as they take on a new shape: \emph{Strong}
might describe an ogre's muscle in one world and a mech suit's hydraulics
in the next. See each genre chapter for local translations.

\subsection*{Special Traits}
Some things are not a bigger pool --- they are a door that dice alone
cannot open. \keyword{Special Traits} are exceptional, often impossible
bodies: wings, a giant's frame, a hide that shrugs off knives. At Body
\rating{3} you choose \dice{1} Special Trait; at Body \rating{4} you
choose \dice{2} in total. Choose from the list at the end of this
chapter.

% ---- Section 2: Rules -----------------------------------------------
\section{Body in Play}

The \keyword{Body} Stat serves double duty: it measures raw physical
power \emph{and} the size of your Body pool --- how much physical harm
you can shrug off before you go down. Mind and Face work the same way
for mental strain and social standing; see \emph{Consequences} in
\emph{How to Play}.

Body rank also determines what physical form you take when you cross
into a new realm --- an exceptional body stays exceptional, but it wears
a different shape. Your Stats, Tags, and Traits all travel with you; only
their fictional dressing changes.

% ---- Section 3: Options ---------------------------------------------
\section{Your Body}

\subsection*{Common Body Tags}
Ten tags, paired into five weight classes so no single pick outshines the
rest --- each pair covers the same kind of ground at roughly the same
strength.

\skillcat{Power}
\begin{itemize}
  \item \atrtag{body}{Strong} --- Lift, break, and carry what ordinary people cannot.
  \item \atrtag{body}{Natural Weapon} --- Claws, fangs, horns, or spines that need no forging.
\end{itemize}

\skillcat{Speed \& Grace}
\begin{itemize}
  \item \atrtag{body}{Fast} --- Outrun, outreact, and strike before others can respond.
  \item \atrtag{body}{Graceful} --- Move with unnatural balance and precision, even where others would stumble.
\end{itemize}

\skillcat{Resilience}
\begin{itemize}
  \item \atrtag{body}{Tough} --- Shrug off wounds, poison, and exhaustion that would fell anyone else.
  \item \atrtag{body}{Regenerating} --- Recover from injury and hardship faster than anyone has a right to.
\end{itemize}

\skillcat{Senses \& Stealth}
\begin{itemize}
  \item \atrtag{body}{Keen Senses} --- See, hear, or smell far beyond human limits, day or night.
  \item \atrtag{body}{Camouflage} --- Your form blends into your surroundings at will.
\end{itemize}

\skillcat{Presence}
\begin{itemize}
  \item \atrtag{body}{Striking} --- Beautiful, magnetic, or impossible to ignore.
  \item \atrtag{body}{Intimidating} --- Your presence alone puts people on edge.
\end{itemize}

\begin{notebox}[Designer's Note]
Resolved 2026-09-22: yes, Body Tags cap at \dice{5} per roll, same as
Magic and everything else --- see \emph{Five Tags, Maximum} in
\emph{How to Play}. One system-wide rule now, instead of a per-column
number to remember.
\end{notebox}

% REAL RULE, 2026-09-26: Body Feats renamed Special Traits and moved here
% from the retired Special Feats chapter, with a new list. Same allotment:
% 1 at rank 3, 2 at rank 4.
\section{Special Traits}
\artinline{feat-body-shapeshifter.png}

A Special Trait never needs a roll. It is simply true of your body, all
the time, with nothing to count or recharge. It never adds or removes
dice; it changes what you can attempt at all. Doing something
\emph{hard} with a Trait is still an ordinary roll: flying is free, but
outflying a dragon is a \keyword{Body} roll. A Trait does exactly what
it says and no more; when a situation is unclear, the GM reads it
narrowly. You may build your own Trait with your GM, as long as it feels
about as strong as the ones below.

\begin{feat}{Burrows}
Dig through earth, sand, and snow as fast as you walk, leaving a tunnel
behind you or letting it collapse. Solid rock and worked stone stop you.
\end{feat}

\begin{feat}{Flying}
Fly under your own power, as fast as you can run, for as long as you
like. Carrying anyone with you is a \keyword{Body} roll.
\end{feat}

\begin{feat}{Fully Animal}
Choose one animal when you take this Trait, anything from a mouse to an
elephant. You \emph{are} that animal, with its body and everything that
body does: a hawk flies, a fish breathes water, a mouse fits through a
crack in the wall. You keep your own mind and voice, but not your
hands.
\end{feat}

\begin{feat}{Giant}
You stand two or three times as tall as an ordinary person, with the
weight to match. You step over walls, wade rivers, and carry what a
cart would. Doorways, beds, boats, and blending into a crowd are now
your problems.
\end{feat}

\begin{feat}{Half Animal}
Choose one animal when you take this Trait. Your body blends with
it --- a centaur, a minotaur, a harpy, a merfolk --- and you keep your
hands and your voice. Choose which of its features you have; each one
does what it does for the animal, so a harpy's wings fly and a
merfolk's gills breathe water.
\end{feat}

\begin{feat}{Mimic}
Reshape your body to look like someone else: your height and weight,
your face, your voice, your hair, your skin, and any marks or scars.
You can pass as another kind of person or copy a specific person you
have studied for a few minutes. You cannot change your size much or
your basic shape, and your clothes and gear do not change with you.
You gain none of their abilities or memories, and acting like them is
still up to you: the GM may call for \keyword{Face} rolls with anyone
who knows them well. Changing takes a moment and lasts until you
change again.
\end{feat}

\begin{feat}{Monstrous Form}
You are something nature never made: a dragon, a demon, a living
statue, a knot of tentacles. Work out with your GM what it is. You have
what that body naturally has, such as scales, horns, stone skin, or a
dozen grasping limbs, and anyone who sees you knows you are no ordinary
creature. To fly, you also need \emph{Flying}; to breathe fire, you
also need \emph{Monstrous Weapon}.
\end{feat}

\begin{feat}{Monstrous Weapon}
Choose one weapon no ordinary creature has: a breath of fire, a
venomous sting, a petrifying gaze, a scream that shatters glass. It is
always ready and never runs out. Using it on someone is an ordinary
roll, but it does what no sword can: set a building alight, freeze a
river, turn a guard to stone.
\end{feat}

\begin{feat}{Robust}
A few hours' sleep a night is all you need. Poison and disease pass
through you harmlessly. You do not need to breathe, so you survive
underwater, in smoke, and in the vacuum of space, and the cold and
pressure of the void and the deep do you no harm. Fire still burns.
\end{feat}

\begin{feat}{Spins Webs}
Spin silk as strong as rope, as much as you like: bridge a chasm, bind a
prisoner, climb your own lines up a cliff. Sticky strands trap anyone
who blunders into them, and breaking free is a \keyword{Body} roll.
Fire burns through it.
\end{feat}

\begin{feat}{Tough Skin}
Simple weapons in the hands of mooks cannot hurt you: clubs, knives,
spears, arrows, a street thug's pistol. You brush them off. Named
foes, heavy weapons, and magic still can.
\end{feat}
